\documentclass[12pt]{article}
\usepackage{ctex}
\usepackage{amsmath, amssymb}
\usepackage{geometry}
\geometry{a4paper, margin=2.5cm}
\usepackage{enumitem}
\title{第 5 章\ 质点的动量矩与动量矩守恒\ 例题}
\author{}
\date{}

\begin{document}
\maketitle

\begin{enumerate}[label=\textbf{例\arabic*.}, leftmargin=*]

% 例 1
\item 一质点 $m$ 位于 $xOy$ 平面内 $\vec{r} = (3\vec{i} + 4\vec{j})\,\mathrm{m}$，速度 $\vec{v} = (2\vec{i} - 5\vec{j})\,\mathrm{m/s}$。求该质点对原点的动量矩。

% 例 2
\item 质量 $m$ 的质点做半径 $R$ 的匀速圆周运动，角速度 $\omega$。求质点对圆心和对圆周上某固定点的动量矩大小。

% 例 3
\item 一质量 $m = 0.1\,\mathrm{kg}$ 的小球沿半径 $R = 0.5\,\mathrm{m}$ 的圆环在水平面内做匀速圆周运动，速率 $v = 2\,\mathrm{m/s}$。求小球对圆心的动量矩大小和方向。

% 例 4
\item 圆锥摆：质量 $m$ 的小球在水平面内做半径为 $r$、角速度为 $\omega$ 的匀速率圆周运动，轻杆与竖直方向夹角为 $\alpha$。求小球对悬挂点 $B$ 的动量矩大小和方向。

% 例 5
\item 一质量 $m$ 的质点做直线运动，$t$ 时刻位置 $x = 2t^3 + t$（m）。求对原点的力矩 $\vec{M}$ 表达式。

% 例 6
\item 人造卫星在椭圆轨道上运行，地球中心可看作固定点。近地点距地面 $439\,\mathrm{km}$，远地点距地面 $2384\,\mathrm{km}$，近地点速度 $8.12\,\mathrm{km/s}$，地球半径 $6370\,\mathrm{km}$。求卫星在远地点的速度。

% 例 7
\item 质量 $m$ 的小球系在绳的一端，绳穿过铅直套管，使小球限制在光滑水平面上运动。先使小球以速度 $v_0$ 绕管心做半径 $r_0$ 的圆周运动，然后向下拉绳，使小球轨迹最后成为半径 $r$ 的圆。求：(1) 小球距管心 $r$ 时速度 $v$ 的大小；(2) 绳从 $r_0$ 缩短到 $r$ 过程中力 $F$ 所做的功。

% 例 8
\item 变半径旋转运动：质点 $m$ 在有心力作用下运动，初始半径 $r_0$、角速度 $\omega_0$。求半径 $r$ 处的角速度。

% 例 9
\item 质量 $m_1$、$m_2$ 的两个小孩各抓着跨过滑轮绳子两端。设滑轮质量和轴上摩擦都可忽略。(1) 若两小孩质量相等，初始都不动，其中一个用力向上爬，问哪一个小孩先到达滑轮？(2) 若 $m_1 \ne m_2$，情况如何？

% 例 10
\item 一宇宙飞船去考察一质量为 $M$、半径为 $R$ 的行星，当飞船静止于空间距行星中心 $4R$ 时，以速度 $v_0$ 发射一质量为 $m$ 的仪器。要使该仪器恰好掠过行星表面，求 $\theta$ 角及着陆滑行（掠过表面时）的初速度。

% 例 11
\item 电子质量 $m$，在原子核库仑引力 $F = -k/r^2$ 下做椭圆轨道运动。证明：电子对原子核的动量矩守恒，且面积速度为常数。

% 例 12
\item 哈雷彗星绕太阳运动的轨道偏心率 $e = 0.967$，近日点距离 $r_1 = 8.8 \times 10^{10}\,\mathrm{m}$，远日点距离 $r_2 = 5.3 \times 10^{12}\,\mathrm{m}$。求：(1) 近日点速率与远日点速率之比；(2) 近日点速率（已知太阳引力常数 $GM_\odot = 1.33\times 10^{20}\,\mathrm{m^3/s^2}$）。

\end{enumerate}

\end{document}
